\begin{claim}{D.2.11}{Proved}{Trivial periodic Lie algebras jointly detect every target}
Fix a prime $p$ and a finite positive height $n$. Let
$\mathcal V=\Sp_{T(n)}$ or $\Sp_{K(n)}$, and put
$\mathfrak L=\operatorname{Lie}(\mathcal V)$.
The trivial Lie algebras generate $\mathfrak L$ under small
colimits and retracts. In fact a countable set of trivial Lie
algebras suffices. In particular,
\[
 \Map_{\mathfrak L}(\operatorname{triv}U,Z)\simeq *
 \quad\hbox{for every }U\in\mathcal V
 \quad\Longrightarrow\quad Z\simeq0.
 \tag{D.2.11a}
\]

For the ordinary-space consequence, use Heuts's fully faithful
realization of the unstable $K(2)$-local category
$\mathcal M_{K(2)}$ inside pointed $p$-local spaces. For
\[
 B=K(\FF_p,2),\qquad A_r=Q\Sigma^rB,\qquad r\geq0,
\]
one has
\[
 Z\in\mathcal M_{K(2)},\quad
 \Map_*(A_r,Z)\simeq *
 \quad\Longrightarrow\quad Z\simeq *.
 \tag{D.2.11b}
\]
Thus no noncontractible target in that unstable periodic category
can separate $B$ from the reversible class generated by $A_r$.
This does not compute the ordinary-space nullification in
\textup{(D.2.9b)}, or compare it with periodic localization.
\end{claim}
\begin{proof}
We first prove the categorical assertion, including its set-size
qualification, and then establish the ordinary mapping-space
interface needed for \textup{(D.2.11b)}.

\medskip\noindent
\textbf{Products and the generated class.}
Corollary~B of \citet{I_BlansHeuts2026} proves that the cartesian
product in $\operatorname{Lie}(\Sp)$ preserves weakly contractible
colimits separately in each variable. This property passes to
$\mathfrak L$. Indeed stable Bousfield localization induces a
reflection on spectral Lie algebras whose underlying spectrum is
the localized underlying spectrum. Its operations are obtained
by localizing the original operations and using
\[
 L_E\bigl(\mathbb L(k)\wedge X^{\wedge k}\bigr)
 \simeq
 L_E\bigl(\mathbb L(k)\wedge(L_EX)^{\wedge k}\bigr).
\]
Here $E=T(n)$ or $K(n)$, and the equivalence follows because
$E$-acyclic spectra form a tensor ideal. Operadic composition and
its coherences localize as well. This is the construction made
explicit for the $K(n)$-local reflection in
\cite[Proposition~4.26]{I_Heuts2021}.
The reflection is left exact: limits of algebras are created on
underlying spectra, and stable localization is exact. Apply it
to the product-colimit comparison in ordinary spectral Lie
algebras. Local colimits are reflected ordinary colimits, and
finite products are preserved, proving the assertion in
$\mathfrak L$.

Suppose for the moment that $\mathcal G$ is a collection of
trivial algebras closed under finite products and containing
zero. Let $\mathcal C$ be its closure under small colimits and
retracts. Then $\mathcal C$ is closed under finite products.
Here is the size-independent closure argument. First fix
$A\in\mathcal G$ and consider the class of $X$ for which
$X\times A\in\mathcal C$. It contains $\mathcal G$, since
products of trivial algebras are trivial algebras on direct sums.
It is closed under retracts. Given a diagram $(X_i)_{i\in I}$
in this class, adjoin an initial vertex to $I$ and put zero
there. The extended diagram $X'$ has the same colimit as the
original diagram, while its indexing category $I^{\triangleleft}$
is weakly contractible. Consequently
\[
 \left(\colim_{i\in I}X_i\right)\times A
 \simeq
 \colim_{i\in I^{\triangleleft}}(X'_i\times A).
 \tag{D.2.11c}
\]
The additional value on the right is $0\times A=A$, which is in
$\mathcal C$, and every other value is in $\mathcal C$ by
assumption. The same argument includes the empty diagram.
Thus $X\times A\in\mathcal C$ for every $X\in\mathcal C$.
Now fix any $Y\in\mathcal C$ and repeat the argument for the
class of $X$ with $X\times Y\in\mathcal C$. It contains all
generators by the first step, and the added initial value is
now $Y\in\mathcal C$. This proves binary-product closure, hence
finite-product closure. We have not asserted preservation of
arbitrary colimits by products: the extra initial value in
\textup{(D.2.11c)} is essential.

\medskip\noindent
\textbf{A fibre sequence as an action-bar realization.}
For a fibre sequence $F\to E\to D$ in $\mathfrak L$ there is an
augmented simplicial object with terms
\[
 [q]\longmapsto F\times(\Omega D)^q
\]
whose realization is $E$. Its face maps encode the canonical
$\Omega D$-action on $F$; they need not describe a trivial action.

To construct and check it, form the \v{C}ech nerve of $0\to D$.
This is the bar construction of the group object $\Omega D$.
Its realization is $D$: the conservative forgetful functor
$U:\mathfrak L\to\mathcal V$ preserves finite limits and
geometric realizations, and in the stable category $\mathcal V$
the corresponding realization is
$\Sigma\Omega U(D)\simeq U(D)$.
Preservation of geometric realizations follows from preservation
of sifted colimits by the spectral Lie monad, whose tensor
products preserve colimits separately. Equivalently, this is
the bar calculation in the proof of
\cite[Proposition~4.11]{I_BlansHeuts2026}.

Pull this augmented nerve back along $E\to D$. Its $q$th term
is $F\times(\Omega D)^q$, by the finite-limit description of
the \v{C}ech action. After applying $U$, realization commutes
with this pullback: finite limits in a stable category are
finite colimits, and geometric realization is exact. The
realized augmentation therefore becomes
\[
 U(E)\mathop{\times}_{U(D)}
       |U(\check C(0\to D))|\simeq U(E).
\]
Conservativity proves the assertion in $\mathfrak L$. In
particular, for the product-closed class constructed above,
\[
 F\in\mathcal C,\quad \Omega D\in\mathcal C
 \quad\Longrightarrow\quad E\in\mathcal C.
 \tag{D.2.11d}
\]
Every simplicial term lies in $\mathcal C$, as does its
realization. The premise concerns $\Omega D$, not merely $D$.
No join axiom or general fibre/base-to-total closure of a
nullification kernel is used.

\medskip\noindent
\textbf{Finite weights and an explicit countable family.}
Write $\mathbb S_{\mathcal V}$ for the local unit. For every
$\ell\in\ZZ$, let
\[
 G_\ell=\operatorname{Free}_{\mathbb L}
                (\Sigma^\ell\mathbb S_{\mathcal V}),
 \qquad J_{\ell,j}=(G_\ell)_j
\]
denote the free algebra and its weight-$j$ spectrum. By
\cite[Theorem~4.22 and Corollary~4.23]{I_Heuts2021},
$J_{\ell,j}=0$ for $j>2p^n$, and already for $j>p^n$ when
$\ell$ is odd. The stated range explicitly includes every
integer $\ell$, including negative integers; the proof treats
negative shifts by partition-complex duality. For $K(n)$-local
spectra apply the monoidal reflection to the $T(n)$-local
formula.

Take $\mathcal G$ to consist of the algebras
\[
 \operatorname{triv}
 \left(\bigoplus_{a=1}^t
       \Sigma^{-k_a}J_{\ell_a,j_a}\right),
 \qquad
 \begin{gathered}
 t\geq0,\quad k_a\geq0,\quad \ell_a\in\ZZ,\\
 1\leq j_a\leq2p^n.
 \end{gathered}
 \tag{D.2.11e}
\]
This is a countable set of specified objects. It contains zero
and is closed under finite products and loops. Let $\mathcal C$
again be its colimit and retract closure. The preceding two
arguments apply to this actual set, without supposing the
existence of a nullification generated by a proper class.

For $m\geq0$, form the positively graded algebra
$(G_\ell)_{\leq m}$ by setting weights above $m$ equal to zero.
Operations add input weights, so these truncations and their
projections are well-defined operadic algebras and algebra maps.
There are fibre sequences
\[
 \operatorname{triv}J_{\ell,m}\longrightarrow
 (G_\ell)_{\leq m}\longrightarrow(G_\ell)_{\leq m-1}
 \qquad(m\geq1).
 \tag{D.2.11f}
\]
The fibre has only weight $m$. Its operations of arity at least
two land in weight at least $2m>m$ and vanish. This verification
is in the category of graded operadic algebras, so it includes
the inherited higher coherences, not just the binary bracket.

Induct on $m$, proving simultaneously that
\[
 \Omega^k(G_\ell)_{\leq m}\in\mathcal C
                   \quad\hbox{for every }k\geq0.
 \tag{D.2.11g}
\]
For $m=0$ the object is zero. Apply $\Omega^k$ to
\textup{(D.2.11f)}. Its fibre is
$\operatorname{triv}(\Sigma^{-k}J_{\ell,m})$, a generator in
\textup{(D.2.11e)}. The loop of its base is
$\Omega^{k+1}(G_\ell)_{\leq m-1}$, which belongs to
$\mathcal C$ by the induction hypothesis for this value of
$k+1$. Formula~\textup{(D.2.11d)} proves the induction step.
At $m=2p^n$ this gives $G_\ell\in\mathcal C$ for every
integer $\ell$. This induction does not assume that
$\mathcal C$ is generally closed under loops.

Local sphere shifts generate $\mathcal V$ under colimits:
present an ordinary spectrum by sphere cells and apply the
localization. The free algebra functor preserves colimits, so
every free Lie algebra lies in $\mathcal C$. Every algebra is
the geometric realization of its free-forgetful monadic bar
resolution. Hence $\mathcal C=\mathfrak L$.

The class of objects $X$ with $\Map_{\mathfrak L}(X,Z)\simeq *$
is closed under colimits and retracts. If it contains all the
trivial algebras, or just the countable family
\textup{(D.2.11e)}, it therefore contains every algebra. Taking
$X=Z$ makes the identity nullhomotopic and gives $Z\simeq0$.
This proves \textup{(D.2.11a)} and the stronger countable
assertion. In particular nullification at that countable family,
if expressed as an accessible localization, is the zero
localization. The argument itself only used the explicit
closure class and mapping-space orthogonality.

\medskip\noindent
\textbf{The precise ordinary-to-periodic interface.}
In the remainder $n=2$. Use the fully faithful models of
$\mathcal S_{v_2}$ and $\mathcal M_{K(2)}$ in
\cite[Definition~3.11, Corollary~2.8 and Proposition~4.26]
{I_Heuts2021}. Let $d_3$ be the connectivity parameter there,
choosing the lower parameter with $d_2\leq d_3$.
Here $L_i^f$ means finite-space nullification at the chosen
type-$(i+1)$ suspension, as in that source.

Suppose a spectrum $H$ is $(d_3+1)$-connective and is both
rationally and $T(1)$-acyclic. Put $X=\Omega^\infty H$.
The identity $\Phi_i\Omega^\infty=L_{T(i)}$, recorded on
page~4 of \citet{I_Heuts2021}, makes its lower periodic
homotopy vanish. Theorem~3.7 of that paper, at height one,
gives $L_1^fX\simeq *$. Proposition~3.6 and the identity
$L_1^fL_2^f=L_1^f$ show that $L_2^fX$ is still
$d_3$-connected and $L_1^f$-acyclic. Thus it belongs to the
specified full subcategory $\mathcal S_{v_2}$ and the
monochromatic functor satisfies $MX\simeq L_2^fX$.

Every realized $Z\in\mathcal M_{K(2)}$ is $L_2^f$-local.
If $F_Z\in\operatorname{Lie}(\Sp_{K(2)})$ is its Lie
algebra, nullification adjunction, fullness, and the
identification of the periodic infinite-loop functor with
the trivial Lie algebra functor give
\[
 \Map_*(\Omega^\infty H,Z)
 \simeq
 \Map_{\operatorname{Lie}(\Sp_{K(2)})}
       (\operatorname{triv}L_{K(2)}H,F_Z).
 \tag{D.2.11h}
\]
For the last identification use Theorem~2.6 and the discussion
on page~8 of \citet{I_Heuts2021}, followed by its
$K(2)$-local reflection. The connective hypotheses justify
the ordinary-space mapping adjunction; we have not treated
$M$ as a reflection from all pointed spaces.

\medskip\noindent
\textbf{Obtaining every trivial source from $A_r$-acyclic spaces.}
Put $C_0=\Sigma^\infty B$ and $E_0=L_{K(2)}C_0$.
The spectrum $C_0$ is rationally and $T(1)$-acyclic. For
the second assertion without an implicit higher telescope
conjecture, first use
\cite[Section~10.11]{EH5Bousfield1997} to get
$K(1)_*B=0$ in reduced homology. Snaith splitting makes
$K(1)_*QB=0$, since each extended power of $C_0$ is
$K(1)$-acyclic. To make the $K(1)$ versus complex $K/p$
interface explicit, Theorem~11.3 makes $QB\to *$ a virtual
$K(1)$-equivalence; immediately before Theorem~11.11,
Bousfield identifies virtual $K(1)$-equivalences with virtual
$K/p$-equivalences. Theorem~11.11 then gives zero
$v_1$-periodic homotopy for the connected $H$-space $QB$
(the same implication summarized in Corollary~11.12).
Finally $\Phi_1QB\simeq L_{T(1)}C_0$ gives the assertion.
It follows that every object of $\loc(C_0)$ has the same
two lower stable acyclicity properties.

We need a unit retract in $\Sp_{K(2)}$. The object $E_0$ is
dualizable by \cite[Corollary~8.8]{I_HoveyStrickland1999}.
The explicit calculation of
\cite[Theorem~11.1(a)]{I_RavenelWilson1980}, with
$n=q=2$ and cyclic group of order $p$, gives unreduced
$K(2)$-rank $p$, concentrated in even degrees. The extension
to $p=2$ is explicitly recorded in
\cite[Remark~4.4]{I_HopkinsRavenelWilson1998}; see also its
Corollary~1.3. Therefore $K(2)_*E_0$ has even reduced rank
$p-1$.

For odd $p$, apply the symmetric graded K\"unneth isomorphism
to the categorical trace of $\id_{E_0}$. It gives multiplication
by $p-1$ on $K(2)$-homology. Thus
\[
 \mathbb S_{K(2)}\xrightarrow{\operatorname{coev}}
 E_0\otimes DE_0\xrightarrow{\operatorname{ev}\circ\tau}
 \mathbb S_{K(2)}
\]
is an equivalence, and rescaling the second map gives a
retraction onto the unit. At $p=2$ no homotopy-commutativity
of Morava $K$ is assumed. The $K(2)$-ranks of $E_0$ and
$DE_0$ are both one. Coevaluation is nonzero on
$K(2)$-homology, since otherwise a duality triangle would
induce the zero map on $K(2)_*E_0$. It is therefore an
isomorphism between rank-one graded modules, and hence an
equivalence of local spectra. This proves in every case that
\[
 \mathbb S_{K(2)}\text{ is a retract of }E_0\otimes DE_0.
 \tag{D.2.11i}
\]
Here $D$ is duality in $\Sp_{K(2)}$ and $\otimes$ its
localized tensor product.

Fix $r\geq0$ and an arbitrary $U\in\Sp_{K(2)}$. In ordinary
spectra form
\[
 V=\tau_{\geq N}(DE_0\otimes U),\qquad H=C_0\wedge V,
 \tag{D.2.11j}
\]
where $N$ is large enough that $H$ is at least
$\max\{r+2,d_3+1\}$-connective. A bounded-above spectrum
is $T(2)$-acyclic: if $T(2)$ is the telescope of a positive
$v_2$ self-map on a finite spectrum $T$, then the terms
$\Sigma^{-kd}(T\wedge W)$ have zero homotopy in any fixed
degree once $k$ is large. The positive self-map is supplied by
\cite[Theorem~9]{HopkinsSmith1998}. Thus their telescope
vanishes for bounded-above $W$, and $W$ is also
$K(2)$-acyclic. Applying this to the complementary
Postnikov section in \textup{(D.2.11j)} gives
\[
 L_{K(2)}V\simeq DE_0\otimes U,\qquad
 L_{K(2)}H\simeq E_0\otimes DE_0\otimes U.
 \tag{D.2.11k}
\]

The spectrum $H$ belongs to $\mathcal I_p=\loc(C_0)$:
this localizing subcategory is a tensor ideal in spectra.
It has $p$-primary integral homology, and its chosen
connectivity allows D.2.9 to be applied. Consequently
$\Omega^\infty H$ is ordinarily $A_r$-acyclic. Its
rational and $T(1)$-acyclicity also permits
\textup{(D.2.11h)}. If $Z$ satisfies the premise of
\textup{(D.2.11b)}, these facts give
\[
 \Map_{\operatorname{Lie}(\Sp_{K(2)})}
 \bigl(\operatorname{triv}(E_0\otimes DE_0\otimes U),F_Z\bigr)
 \simeq *.
\]
By \textup{(D.2.11i)}, $\operatorname{triv}U$ is a retract
of this source. Hence $F_Z$ is orthogonal to every trivial
Lie algebra. The already proved categorical assertion makes
$F_Z=0$, and fullness of the realization gives $Z\simeq *$.
The auxiliary $N$ can be increased for every fixed $r$, so
this proves \textup{(D.2.11b)} for all $r\geq0$, including
those below the periodic model's connectivity cutoff.

The proof covers every periodic target, even those with an
unbounded Lie weight filtration. The finite filtration was
used only for a set of free generators of the category.
It does not imply that ordinary $A_r$-nullity is preserved
by periodic localization, and therefore does not settle
\textup{(D.2.9b)}.
\end{proof}
